%------------------------------------------------------------------------------------
%  SECTION 2.1 : DÉMONSTRATION DE TOUS LES NIVEAUX DE SECTIONNEMENT
%------------------------------------------------------------------------------------
\section{Hiérarchie Intégrale des Titres et Divisions Structurelles}
\label{sec:niveaux_sectionnement}

Un document scientifique clair et pédagogique doit s'appuyer sur une structure hiérarchique sans ambiguïté.
La classe \texttt{ClasseGuillaume}, construite sur le socle KOMA-Script (\texttt{scrbook}), personnalise l'ensemble des niveaux de sectionnement du document, depuis la grande division générale jusqu'au paragraphe en ligne.

\subsection{Le niveau macroscopique : Les Parties (\texttt{\textbackslash part})}
\label{subsec:parties}

La commande \texttt{\textbackslash part\{Titre de la Partie\}} regroupe plusieurs chapitres au sein d'une unité thématique majeure.
Dans ce modèle, chaque partie génère une page de garde majestueuse, centrée et aérée, portant un numéro romain.
Les parties permettent d'architecturer les documents de plus de 60 pages sans saturer la numérotation des chapitres.

\subsection{Le niveau capitulaire : Les Chapitres numérotés et non numérotés}
\label{subsec:chapitres}

Les chapitres constituent le cœur du découpage :
\begin{itemize}
  \item \textbf{Chapitre numéroté (\texttt{\textbackslash chapter\{...\}}) :}
  La classe personnalise le titre en adoptant la police {\AsapFont Asap Condensed}, la couleur primaire du document, et une numérotation en chiffres romains (\textsc{Chapitre I}, \textsc{Chapitre II}).
  Le titre est harmonieusement centré et séparé du texte par des espacements verticaux soigneusement calibrés (\texttt{beforeskip = -10pt}, \texttt{afterskip = 20pt}).
  \item \textbf{Chapitre non numéroté KOMA-Script (\texttt{\textbackslash addchap\{...\}}) :}
  En \LaTeX{} standard, la commande \texttt{\textbackslash chapter*\{...\}} présente un défaut majeur : elle n'insère pas le chapitre dans la table des matières et ne met pas à jour les en-têtes de page.
  La commande KOMA-Script \texttt{\textbackslash addchap\{...\}} résout parfaitement ce problème : elle produit un chapitre sans numéro, mais l'enregistre automatiquement dans la table des matières et actualise les en-têtes (tel que réalisé pour le Résumé et les Remerciements).
\end{itemize}

\subsection{Les divisions intermédiaires : Sections et Sous-sections}
\label{subsec:sections_subsections}

La numérotation des sections adopte une convention hiérarchique arabe préfixée par le numéro du chapitre courant :
\begin{itemize}
  \item \textbf{Section (\texttt{\textbackslash section\{...\}}) :}
  Numérotée par exemple \textbf{II.1}, mise en valeur par la couleur secondaire (\texttt{secondaryColor}) et un corps de texte large en gras.
  \item \textbf{Section non numérotée dans la TOC (\texttt{\textbackslash addsec\{...\}}) :}
  À l'instar de \texttt{\textbackslash addchap}, elle insère un en-tête de section sans numéro tout en figurant dans la table des matières.
  \item \textbf{Sous-section (\texttt{\textbackslash subsection\{...\}}) :}
  Numérotée par exemple \textbf{II.1.1}, typographiée en gris foncé soutenu (\texttt{darkgray}) pour créer un contraste hiérarchique doux.
  \item \textbf{Sous-sous-section (\texttt{\textbackslash subsubsection\{...\}}) :}
  Offre un niveau de granularité supplémentaire pour les démonstrations techniques pointues ou les descriptions expérimentales détaillées.
\end{itemize}

Voici ci-dessous une démonstration immédiate d'une sous-sous-section :

\subsubsection{Exemple pratique de sous-sous-section technique}
Cette sous-sous-section permet de délimiter un point précis d'analyse sans surcharger la numérotation globale.
Elle s'intègre harmonieusement dans le flux de lecture sans créer de rupture visuelle excessive.

\subsection{Les divisions terminales : Paragraphes avec puce en ligne}
\label{subsec:paragraphes_lignes}

En \LaTeX{} standard, l'environnement \texttt{\textbackslash paragraph\{...\}} engendre fréquemment une mise en page gauche : soit il renvoie maladroitement à la ligne, soit il s'accole de manière peu esthétique.

\begin{boiteSynthese}{L'innovation de ClasseGuillaume pour \texttt{\textbackslash paragraph}}
  La classe redéfinit entièrement la commande \texttt{\textbackslash paragraph} :
  \begin{verbatim}
  \RedeclareSectionCommand[
    indent = 5mm,
    afterskip = -1em,
    font = {\color{primaryColor}\bfseries\textbullet~}
  ]{paragraph}
  \end{verbatim}
  Le titre du paragraphe est automatiquement précédé d'une \textbf{puce discrète colorée} (\texttt{\textbullet}), indenté de \qty{5}{\milli\meter}, et le texte du corps s'enchaîne \textbf{directement sur la même ligne} après une espace fine.
\end{boiteSynthese}

Démonstration en situation réelle :

\paragraph{Critère de modularité}
Le code source est segmenté en unités logiques autonomes pour autoriser une maintenance aisée et des révisions collaboratives sans friction.

\paragraph{Critère d'accessibilité typographique}
Les contrastes de couleurs, la lisibilité des polices et l'aération spatiale garantissent un confort de lecture optimal pour les jurys et relecteurs.

\subparagraph{Niveau sub-paragraph}
Pour les nomenclatures ou sous-critères nécessitant une hiérarchisation ultime, la commande \texttt{\textbackslash subparagraph} offre une indentation spécifique tout en conservant une discrétion absolue.

